% =====================================================================
\section{Experiments}\label{sec:experiments}
% =====================================================================

This section evaluates the system in the order the pipeline builds it. \Cref{sec:dataset} sets out
the shared setup---the human-annotated flaw benchmark and the recall metrics and matching protocol
every result is scored under. \Cref{sec:gen-experiments} takes the \emph{physics-guided generator}:
whether the geometric control governs the motion, when the trajectory commits, and how far the scene
library reaches. \Cref{sec:eval} takes the \emph{\PhyReAct{} critic}: its flaw recall on the
evaluation core, against the annotated reference and published evaluators, and its agreement with
human ratings. \Cref{sec:icl-exp} takes the \emph{context-training self-evolution} of the
critic---whether lessons distilled from recorded misses raise recall without touching a
weight---and \cref{sec:refine} closes the critic's outward channel, feeding a verdict back as a
prompt rewrite that regenerates the clip. The evaluations use distinct held-out settings that should
not be conflated; \cref{sec:eval-scope} sets them side by side.

\subsection{Experimental setup}\label{sec:dataset}
This setup---models, benchmark, and matching protocol---is shared by both systems; each system's own configuration accompanies its results.

\paragraph{Models and specialists.} The critic runs one served vision-language model for both its text-only planner and its video-aware semantic verifier---Qwen3-VL-30B-A3B-Instruct~\cite{qwen3vl2025}, in BF16 without quantization---together with a frozen stack of low-level specialists, one per capability: SAM~3~\cite{sam3_2025}, TAPNext++~\cite{tapnext2025}, a depth model, PaddleOCR, and FlexSED~\cite{flexsed2025}, with a closed model as the fallback verifier. \Cref{tab:specialists} lists, for each capability, the backend that runs it and the object it returns, together with the exact checkpoint, size, and license; every component is frozen at the listed revision for all reported results.

% Merged specialist manifest, three columns. The capability and the frozen model that runs
% it are folded into one "Capability (model)" column, so each capability is a single
% independent row (no \multirow) and the two description columns get the width to stay on
% one line. License column dropped; the one non-commercial license (depth ordering,
% CC-BY-NC-4.0) is stated in the caption. Hand-maintained.
% Tracks come from TAPNext++ (per Sec. 5); SAM 3 does grounding/masks/count.

\begin{table}[t]
\centering\small
\setlength{\tabcolsep}{5pt}
\renewcommand{\arraystretch}{1.25}
\begin{tabularx}{\textwidth}{@{}p{6.35cm} X X@{}}
\toprule
\textbf{Capability (model)} & \textbf{Running backend} & \textbf{Returned object} \\
\midrule
Planner (Qwen3-VL-30B-A3B, 31\,B)~\cite{qwen3vl2025}
  & prompt compiled pre-frame & claims and a typed plan \\
Semantic predicate (Qwen3-VL)
  & dense timestamped pass & learned base state, rationale \\
Event timing (Qwen3-VL)
  & over timestamped frames & interval, absence, or refusal \\
Fallback verify (gpt-5.6-sol)
  & on a claim the dense pass passed & three-valued state, both kept \\
Grounding, masks (SAM~3, 0.9\,B)~\cite{sam3_2025}
  & SAM~3 Video, persistent IDs & per-frame masks, tracked IDs \\
Object count (SAM~3)
  & cardinality of SAM~3 IDs & count, per-frame series \\
Physical measurement (TAPNext++)~\cite{tapnext2025}
  & 11 typed operations over tracks & typed measurement, evidence \\
Depth ordering (SenseNova-7B-MoT)
  & depth model on localized evidence & behind / in-front ordering \\
Camera cancellation (from tracks)
  & dense optical flow subtracted & subject-relative displacement \\
Rendered text (PaddleOCR PP-OCRv5)
  & OCR head + EasyOCR fallback & literal glyph transcription \\
Audio (FlexSED)~\cite{flexsed2025}
  & sound-event detection & sound intervals on the track \\
\bottomrule
\end{tabularx}
\caption{Execution backends and the frozen specialist behind each capability, for the
evaluated run: each row names a capability with the model or checkpoint that runs it (with
its parameter count where published), the backend it runs, and the object it returns.
SAM~3 supplies grounding, masks, and counting, while TAPNext++ supplies the tracks the
eleven typed operations (\cref{sec:physics}) are taken over; the two learned components
(semantic predicate and depth ordering) and the fallback verifier are the only backends
that may resolve a claim, and no holistic video-quality model produces the clip decision.
Every component is frozen at the listed revision for all reported results, and
\texttt{gpt-5.6-sol} is a closed model (SenseNova SOL) served through an internal proxy.
All released components are under permissive licenses (Apache-2.0, MIT, or Meta's SAM
license) except the depth-ordering model, which is CC-BY-NC-4.0 (non-commercial).}
\label{tab:specialists}
\end{table}


\subsubsection{Benchmark}\label{sec:bench-spec}
The benchmark is a corpus of $1{,}500$ AI-generated clips: a person watched each against its prompt and recorded every place the video failed to deliver what the text asked. The generating model is not recorded, no clip is attributed to a specific system, and the human-written fields are ground truth throughout. Each flaw is a localized record with four always-present fields---the \textbf{violated fragment}, a free-text \textbf{rationale}, a \textbf{severity}, and the annotator's \textbf{confidence}---and two optional, a \textbf{time span} and a \textbf{box track}. Two grouping tags---a \emph{modality} tag (seen, heard, or both) and a ten-item \emph{type} tag---are machine-assigned by majority of three zero-temperature model votes and never treated as ground truth. \cref{sec:app-benchmark} gives per-field coverage, the per-type breakdown, and timing distributions.

Three recorded stages (\cref{fig:benchcomp} in \cref{sec:app-benchmark}) narrow the $1{,}500$ clips and $2{,}582$ flaws to the evaluation set: \emph{clip selection} keeps retrievable, dynamic clips; \emph{flaw selection} keeps only ``seen''- or ``both''-tagged flaws a pixel-only critic could find, leaving $1{,}107$. A deterministic seeded procedure, stratified by severity, flaw count, and duration, draws \emph{the frozen core}---$150$ clips and $306$ flaws (a $29$-clip continuity stratum exempt from replacement); a separate pool of $711$ clips is held disjoint from both frozen sets for learning. The core is an optimization set, so results on it measure fit to that core. One clip produced no measurement, so recall is reported over $149$ clips and their $304$ flaw entries, never the full $306$---the two excluded rather than counted as missed, making the reported rate very slightly optimistic.

Unlike clip-level resources\label{sec:bench-vs}---Physion-Eval~\cite{physioneval2026}, VideoPhy-2~\cite{videophy2_2025}, and per-dimension suites such as VBench~\cite{vbench2024} and T2V-CompBench~\cite{t2vcompbench2024}---our corpus ties each failure to the exact prompt words with an optional time span and box track, so an allegation can be matched to a particular annotated failure rather than only correlated with a clip-level score. This localization does not improve reliability: our corpus is single-annotator and contains only flawed clips, whereas VideoPhy-2 provides redundant judgments and majority resolution.
\subsubsection{Reference standard and matching protocol}\label{sec:eval-protocol}
The evaluation core is the reference standard: its $149$ clips carry $304$ annotated flaws---the denominator for all core recall analyses. Because every clip carries at least one flaw, recall is directly measurable but precision is not (a detector flagging everything scores perfect recall), so we report recall alongside the rate at which accusations land on a real flaw. Recall counts flaw entries, not distinct spans, since matching claims each entry at most once. Matching\label{sec:matching}, run after critic inference, cannot affect the critic's outputs. The comparison unit is a sentence pair (annotator vs.\ critic), judged by meaning as a \emph{whole} match (same defect), a \emph{partial} match (critic identifies at least one aspect of a compound annotated defect), or a non-match. Whole and partial pairings are candidate edges for maximum bipartite matching, each finding and flaw used at most once; every flaw is then reported as whole, partial, or missed. The matching judge, distinct from all evaluated components, votes at temperature zero with a retry-and-refallback rule and aborts rather than reporting a number if more than $2\%$ of flaws yield no usable outcome.
\subsection{The physics-guided generator}\label{sec:gen-experiments}

Three questions organize the generator experiments: whether the geometric control overrides
prompt-implied motion (\cref{sec:gen-prompt}), when a conditioned trajectory becomes insensitive to
control withdrawal (\cref{sec:gen-release}), and how the pipeline behaves beyond single-body probes
(\cref{sec:gen-breadth}). The first and third use the deployed generator; the release analysis is a
separate mechanism probe on the second backbone, its withdrawal operation defined in
\cref{sec:gen-withdrawal}.

\subsubsection{Setup: from simulation to a controllable clip}\label{sec:gen-setup}

The deployed backbone is Wan~2.2-VACE-Fun-14B (\code{WanVACEPipeline}, fixed negative prompt, no
appearance reference, deterministic sampling from the recorded seed). Scenes are simulated in MuJoCo
and rendered to $81$ frames at $16$\,fps; as the simulated duration varies, this render clock
stretches playback by $1.8\times$--$2.3\times$, so trajectory shape and event order come from the
simulation but absolute speed does not equal simulator time. Generation is staged---controls
rendered first, a 1.3B checkpoint screening each scene once, visibly degenerate outputs not queued
for the 14B model (a pipeline rule, not a generator property). Surviving scenes are generated in
bf16 at $848\times480$ with UniPC at $50$ steps and guidance $5.0$, conditioned on the control video
(globally normalized MuJoCo depth, or a silhouette for predeclared flat-scene templates), an
all-white mask, and the prompt; the conditioning scale is $s=1.0$ deployed, examined separately at
$\{0.5,0.8,1.0\}$.

\subsubsection{Control fidelity under prompt variation}\label{sec:gen-prompt}

We fix the geometric control---one projectile-depth video---and vary only the motion clause across three prompts sharing subject, scene, camera, material and style: a strong horizontal traverse, a neutral one, or one adversarial to the control (drops nearly straight down, almost no sideways motion). Each runs under three regimes---no control, control over the first $10$ of $50$ steps, control throughout---at three seeds (partial), two (unconditioned), one (full). We track horizontal travel (fraction of frame width) in the finished clip; metric and decision rules were registered before generation.

\begin{figure}[t]
  \centering
  \includegraphics[width=\linewidth]{fig_prompt_trajectories.pdf}
  \caption{Tracked ball path in nine clips at one seed, against the same depth control throughout.
  Each box is a single $848\times480$ frame, $x$ left to right and $y$ downward; each dot is one
  frame, sized by the tracked region's area and shaded from first frame to last. The pale band is
  the control video's ball, in frame for $63$ of its $81$ frames; the number in each box is that
  clip's sideways travel. The three prompts differ only in the closing sentence beside each row;
  the red row asks for the motion the control does not perform. Without the control the path
  follows the clause and the framing shifts with it; with the control throughout, all three
  clauses reproduce its arc; with the control over the first ten steps the path holds the arc and
  departs at the landing.}
  \label{fig:prompttraj}
\end{figure}

The two inputs cross over (\cref{fig:prompttraj}). Without a control the prompt sets the motion: the clause-attributable range is $0.875$ against $0.252$ between seeds. The unconditioned clips also differ in framing---median tracked ball area $8{,}846$ pixels against $211$ under the control---so travel-fraction is not comparable across regimes, and within-regime comparisons carry the argument. With the control throughout, the clause-attributable range falls to $0.001$; in between, with the control over the first fifth of sampling then withdrawn, the two sources swap: clause range $0.060$ against a seed range of $0.427$.

The trajectories locate the shortfall (\cref{fig:prompttraj}). A clip conditioned over the first ten steps tracks the control's arc through the flight and departs only at the landing, stopping near $0.77$ of the frame while the control's ball rolls on to $0.99$: the lost travel is the post-contact roll, not a different flight---the same phase \cref{sec:gen-release} finds most sensitive to early withdrawal. Of the three clauses at this seed, two follow the control from the first frames and the neutral one does not, staying a mean $0.243$ of the frame from the control's path over the first ten frames. The control-set trajectory thus persists through the text-only steps.
\subsubsection{Trajectory commitment under control withdrawal}\label{sec:gen-release}

Moving the withdrawal point locates the noise level at which the simulated trajectory becomes fixed. The withdrawal operation, its second backbone, $\sigma_{\mathrm{rel}}$, and the two variants clearing or keeping the solver's cached derivatives are defined in \cref{sec:gen-withdrawal} (\cref{eq:flow-sigma}). Scenes are one moving rigid body with a closed-form trajectory (one of three yields only five post-contact frames, so its comparisons rest on two), one schedule setting per seed. Since noise falls during sampling, a \emph{high} $\sigma_{\mathrm{rel}}$ withdraws \emph{early}, leaving later steps unconditioned; a clip still following the simulation afterward was fixed before removal, and the earliest such withdrawal marks the commitment point. Three runs probe seed variability ($16$ seeds, one point), transition shape (six points, one seed), and measurement dependence (a denser sweep re-scored with an independent tracker).

\paragraph{Metrics.} Each run pairs against a $\sigma_{\mathrm{rel}}=0$ run sharing simulation, control, prompt, sampler, and seed, both tracked frame by frame (by colour, except the third run's segmentation tracker). \emph{Retention}, the per-axis Pearson correlation between generated and control paths (\cref{sec:app-retention}), is scored separately over \emph{airborne} and \emph{after-contact} windows split at the ground-contact frame $t_c$; near $1$ means the object moved as simulated, and the split localises disagreements a whole-clip value hides. Two prefixed failure counts accompany it: \emph{whole-clip disagreement} (full-clip horizontal $<0.7$) and \emph{late-window disagreement} (whole-clip horizontal $<0.7$ while airborne horizontal $>0.9$).

\paragraph{Airborne and after-contact agreement under an early release.} The first run fixes release at $\sigma = 0.952$ over two scenes, $16$ seeds each, both variants, against the never-released baseline: $96$ clips, $16$ per cell. Releasing at $\sigma = 0.952$ leaves the thrown ball's airborne flight intact---per-seed horizontal agreement drops $0.003$ versus never releasing ($95\%$ CI $0.001$ to $0.004$, $n=16$ paired seeds)---while whole-clip horizontal falls $0.13$ ($95\%$ CI $-0.02$ to $0.28$) and vertical $0.18$ ($95\%$ CI $0.07$ to $0.29$). The disturbance is confined to post-contact motion, where horizontal runs \emph{opposite} to the simulation on $9$ of $16$ seeds at $\sigma = 0.952$ versus $0$ of $16$ when never released. Whole-clip disagreements number $3$ of $16$ against $1$ of $16$, late-window $2$ of $16$ against $0$ of $16$.

\begin{figure}[t]
  \centering
  \includegraphics[width=\linewidth]{fig_injection_sweep.pdf}
  \caption{The release-point sweep: one clip per point, three scenes (including a
  ball rolling off a ramp), scored separately over the airborne window and the
  window after the simulated object first reaches the ground. The horizontal axis
  is the noise level at which the control was switched off, high noise on the left.
  Colour marks the window (blue airborne, rose after contact) and line style the axis
  (solid across-frame, dashed up/down); retention is the correlation between the generated object's tracked path
  and the simulation's, defined in the text. The thrown ball's after-contact
  agreement falls below the threshold in both directions at the earliest withdrawal,
  and the vertical component is the later of the two to recover. Guides mark the
  commitment band and the backbone's expert switch; the grey band marks the region
  below the $0.7$ threshold fixed before scoring.}
  \label{fig:injsweep}
\end{figure}

Sweeping six points on one seed across three scenes traces retention against noise level (\cref{fig:injsweep}). For the thrown ball, $\sigma=0.952$ gives after-contact horizontal $-0.63$ (opposite to the simulation) with airborne still $1.00$; $\sigma=0.903$ or later restores after-contact to $0.98{+}$ through the latest release, $\sigma=0.555$. A dropped ball stays above threshold everywhere: after-contact horizontal is $1.00$ throughout, and vertical---the informative axis---is lowest at $0.78$ ($\sigma=0.952$), rising to $0.99{+}$ from $\sigma=0.903$. A ball off a ramp stays at $0.99{+}$ in both axes. A third run (three thrown-ball points, three seeds, two guidance settings, $18$ clips, segmentation tracker) reproduces the direction: airborne horizontal $1.00$ (mean over $6$ clips at $\sigma=0.952$, $95\%$ CI $0.998$--$1.000$) while after-contact horizontal averages $0.69$ over an interval spanning zero ($-0.04$ to $1.43$, $n=6$), recovering to $0.99$ and $1.00$ at $\sigma=0.903$ and $0.833$; at the latest point the tracker finds no object on $3$ of $6$ clips (low-contrast ball on grass), excluded, widening that interval.

Across runs the trajectory stops responding to control removal in a narrow band, $\sigma \approx 0.94$--$0.96$. Under \cref{eq:flow-sigma} at $S=40$, $\kappa=5.0$ this sits in the high-noise portion of sampling, consistent with the critical-window account~\cite{phasetrans2024,dynregimes2024,critwindows2024} and above the backbone's expert switch at timestep $0.875\,N$~\cite{wan22dual2026}, so motion is fixed before the handoff; we report noise level not step index because changing $S$ or $\kappa$ moves the step but not the level~\cite{sdedit2022,p2train2022}. The two components do not fix together: on the thrown ball, $\sigma=0.9522$ drops after-contact agreement below threshold in both coordinates ($-0.632$ horizontal, $+0.016$ vertical); one point later ($\sigma=0.9025$) horizontal recovers to $+0.981$ while vertical is still $+0.260$, not reaching $+0.894$ until $\sigma=0.8331$; the dropped ball shows the same ordering, so a schedule set by horizontal alone would be too early. (Whole-clip the ordering reverses---vertical $+0.986$ against horizontal $+0.421$---as the large flight excursion dominates the correlation and hides the rebound.)

The failure is specific: flight is reproduced, post-contact behaviour is not. Contacts are a documented failure mode of video generators~\cite{videophy2024}, traced to a weak contact prior; where the prior is weak the control must stay present to override it, so applied throughout it reproduces the contact on every seed and the reversal appears only under early withdrawal. Clearing the cached derivatives changes the whole-clip result by at most $0.04$ horizontally ($95\%$ CI $-0.04$ to $0.13$, $n=16$), the largest per-seed difference on any window or direction.

\FloatBarrier
\subsubsection{Scene-library breadth and failure boundaries}\label{sec:gen-breadth}

The same authoring, simulation and conditioning path scales beyond the single-body probes, producing $1{,}314$ scenes across $66$ \emph{kinds} (a kind is a template, a scene one parameter setting, simulated and generated once); $591$ of $652$ simulations passed a pre-render inspection of simulator output, not generated clips, judged to show the intended behaviour clearly enough to generate from. \Cref{fig:scenebreadth} shows twelve kinds as generated frames at body counts up to $34$, and \cref{fig:scenepairs} three as depth control and generated clip at matched timestamps.

Fidelity does not extend uniformly. Assessed by manual inspection with the colour tracker of \cref{sec:gen-release}, single-object scenes reproduce the simulated trajectory throughout (\cref{fig:gencases}), but a multi-object interaction such as a billiards break does not---neither the deployed backbone nor the alternates produced a collision in which every body moved as specified---and its commitment point is unmeasured, the colour tracker being unable to follow a scatter as one path. The same rigid-body engine produces every scene, deformable and granular kinds built from its particle-and-link primitives (springs, linked segments, contacting bodies, buoyancy), so the library excludes fluid, smoke and finite-element soft-body solves.

\begin{figure}[tb]
  \centering
  \includegraphics[width=\linewidth]{fig_generation_cases.pdf}
  \caption{The same mechanism across two scenes with very different requested
  appearance---an indoor studio and an outdoor field---each shown as its depth
  control above and the resulting generated clip below. Timestamps are per scene,
  chosen from the turning points of its tracked vertical
  path, so each column falls within a moving interval. Appearance changes completely with
  the prompt while the motion stays the simulation's: for the thrown ball the
  generated and control paths agree to a correlation of $1.0000$ horizontally
  and $0.9998$ vertically, with a median disagreement of $0.7$ pixels across and
  $1.1$ pixels down, over the $63$ frames in which the ball is measurable in
  both before it leaves the frame.}
  \label{fig:gencases}
\end{figure}

\begin{figure}[tb]
  \centering
  \includegraphics[width=\linewidth]{fig_scene_breadth.pdf}
  \caption{Twelve of the $66$ scene kinds, one generated frame each. Each tile is
  selected by rule: frame-to-frame change is accumulated over the clip and
  the tile is the frame by which half the total has occurred, which falls during the scene's principal motion rather than its start. The second caption line names the physics
  primitive, and every body count in it is the scene's own parameter. Each clip's simulation passed the
  inspection described in the text, and the tiles were additionally inspected manually.}
  \label{fig:scenebreadth}
\end{figure}

\begin{figure}[tb]
  \centering
  \includegraphics[width=\linewidth]{fig_scene_control_pairs.pdf}
  \caption{Three deformable and granular scenes, each with simulated depth
  control above and generated clip below, at matched timestamps. Columns are placed at equal steps of the clip's
  accumulated change, so none falls in a temporally static interval. Pairs are shown only for scenes
  whose depth control remains distinguishable from its background in every column;
  an object that settles onto the floor takes the same grey as the floor and leaves
  the control blank.}
  \label{fig:scenepairs}
\end{figure}

\subsection{The \PhyReAct{} critic}\label{sec:eval}
The benchmark and metrics are in \cref{sec:dataset}; this subsection reports the critic's recall on the evaluation core, broken down by clip property and specialist trace. How the settings here read together is deferred to \cref{sec:eval-scope}.

\paragraph{Implementation and cost.}
The served Qwen3-VL model (\cref{sec:dataset}) runs with tensor parallelism across two GPUs. Planner and verifier are separate calls: the planner sees the prompt and claim state but no frames, the verifier sees densely decoded, timestamped frames and one declared predicate. Every measurement uses the logged dense-input path with the frozen target, BF16 inference, and pinned tool revisions. Since a detector, a tracker, and a full vision-language judgement differ in cost, cost is reported as model and tool calls per clip rather than wall-clock time, and uncertainty by resampling whole clips, since two flaws in one clip are not independent.

\paragraph{Planner validity.}\label{sec:planner-gaps}
Inspected before end-to-end evaluation without video or tool access, the prompt-only planner produced a valid, compilable plan for all $233$ claims of $22$ annotated clips (one or two calls per clip); $27$ claims ($11.6\%$) carried a specialist action surviving validation, against $2.3\%$ under an earlier tabular format. On diagnostic sets requiring quantitative measurement the routing rate was $19/30$ across counting, spatial, and ordering claims and $10/10$ for synthetic counting prompts---tests of routing and syntax, not verdict accuracy. The supported vocabulary lacks an ``on top of'' predicate and rejects negated measurements; such claims fall back to the generalist verifier.
\subsubsection{Flaw recall on the evaluation core}
\label{sec:headline}
\paragraph{Headline recall.}
Over the $149$ core clips and their $304$ human-annotated flaws, \PhyReAct{} finds \textbf{228
($75.0\%$)} under \cref{sec:matching}---$191$ whole, $37$ partial---and finds every annotated flaw
on $91$ of the $149$ clips. Two components drive matches: the fallback verifier gpt-5.6-sol overrode
a dense-verifier pass to contradiction on $83$ of the $228$, and the on-screen-text recognizer
transcribes malformed glyphs literally rather than normalizing them, so string comparison catches
garbled text.

\paragraph{Versus published evaluators.}\label{sec:external}
We compare \PhyReAct{} against defect-localizing evaluators---question-decomposition and scene-graph methods (per-question answers linkable to claims) and free-text evaluators (unstructured flaw lists)---and test whether CLIPScore tracks aggregate clip properties. Modular video QA is context-only, not evaluated; clip-level agreement with human ratings is deferred to \cref{tab:vs2}.

\begin{wrapfigure}{r}{0.46\textwidth}
  \centering
  \vspace{-10pt}
  \includegraphics[width=0.44\textwidth]{fig_external.pdf}
  \caption{Localized flaw recall vs.\ cost per evaluator, over the same $149$
  clips and $304$ flaws (\cref{sec:matching}). Horizontal axis: mean model and tool calls
  per clip; vertical: recall. Up-and-left is better. Full comparison: \cref{tab:external}
  in \cref{sec:app-external}.}
  \label{fig:external}
  \vspace{-8pt}
\end{wrapfigure}

The same served model---single-pass over the prompt and every frame, no planning, tools, or per-claim verification---matches $222$ of the $304$ flaws at one call per clip; \cref{fig:external} reports it separately, since free-text lists are prompt-sensitive with no per-verdict evidence.

Question decomposition matches \textbf{164 of the $304$ flaws ($53.9\%$)}---$125$ whole, $39$ partial---versus \PhyReAct{}'s $228$, on the same $149$ densely-read clips and model (identical claim lists on $148$). It asks one closed question per claim and accepts the answer, whereas \PhyReAct{} plans measurements and may run several checks per claim. Higher recall is not a higher per-finding match rate: $29\%$ of its $557$ findings match a human flaw versus $26\%$ of \PhyReAct{}'s $871$. About a third of its questions only check object presence (\emph{Is there a vending machine?}), answerable from one frame; twelve clips get no finding. Clip by clip it finds fewer than \PhyReAct{} on $70$, equal on $63$, more on $16$.

As a clip-level diagnostic, CLIPScore correlates across the $149$ clips with severity ($r=+0.035$, $p=0.67$) and flaw count ($r=-0.024$, $p=0.77$); with $86$ clips carrying multiple flaws ($241$ of $304$), a single scalar cannot localize them.
\paragraph{Recall by category and clip property.}\label{sec:bycat}
% GENERATED by analysis/emit_tex.py — do not edit by hand.
% Every number is recomputed from per-item rows; see the analysis scripts.

\begin{table}[t]
\centering\small
\setlength{\tabcolsep}{6pt}\renewcommand{\arraystretch}{1.15}
\begin{tabular}{@{}p{4.4cm} c c c c c@{}}
\toprule
\textbf{What the annotator complained about} & \textbf{Flaws} & \textbf{Whole}
 & \textbf{Part} & \textbf{Found} & \textbf{Rate} \\
\midrule
action or event occurs & 137 & 100 & 14 & 114 & 83\% \\
how many & 37 & 10 & 8 & 18 & 49\% \\
object identity & 30 & 14 & 6 & 20 & 67\% \\
on-screen text & 25 & 21 & 0 & 21 & 84\% \\
spatial relation & 18 & 10 & 1 & 11 & 61\% \\
order and timing & 17 & 10 & 3 & 13 & 76\% \\
attribute & 16 & 10 & 2 & 12 & 75\% \\
camera and style & 10 & 7 & 2 & 9 & 90\% \\
other & 7 & 5 & 1 & 6 & 86\% \\
physical motion & 7 & 4 & 0 & 4 & 57\% \\
\midrule
\textbf{All} & \textbf{304} & 191 & 37 & \textbf{228} & \textbf{75.0\%} \\
\bottomrule
\end{tabular}
\caption{Flaw recall on the evaluation core, by the annotators' own category: 228 of 304
human-annotated flaws over $149$ clips. ``Whole'' counts flaws a single finding covers completely;
``Part'' counts flaws our findings address in aspect without any one being the whole defect; their
sum is ``Found''. Categories are the annotators' labels, assigned before any critic ran. Every
figure is recomputed from one row per (clip, flaw) parsed from the run's own record, which the
parser cross-checks against that record's stated totals before any table is built.}
\label{tab:headline}
\end{table}

By the annotators' own categories (\cref{tab:headline}), recall runs from $90\%$ to $49\%$. Camera
and style ($90\%$) and on-screen text ($84\%$) rank highest, both measured directly against an
explicit target; counting is lowest at $49\%$ because its path reduces a per-frame series to its mode
(\cref{sec:contract}), so an intermittently detected subject reads as absent. Physical motion holds
seven flaws, four found; most concern deformation---cloth folding wrongly, a body bending where it
should not---which lies outside the trajectory measurements, and its recall is read descriptively
given the size. Grouping the $304$ flaws instead by severity, duration, flaw count, and difficulty
stratum keeps subgroup recall within about ten points of the overall rate everywhere---severity~$5$
lowest at $18$ of $28$---so neither severity nor stratum clearly predicts whether the critic finds a
defect (\cref{tab:robust} in \cref{sec:app-robust}).

\subsubsection{Specialist measurement traces}
\label{sec:cases}
\Cref{fig:criticcases} shows six matched findings from an earlier physics-routing pass over the evaluation core; they are not from the run of \cref{tab:headline} and do not contribute to its recall. Five rest on numeric measurements, the text case on a recognition-head transcript compared as a string. The selection rule was fixed in advance---clips whose every annotated flaw was linked to a finding, and only where the raw per-check record attributed that finding to a specialist measurement and stored the values---excluding ten candidates for mismatched linkage, a verifier-produced finding, or, once, an invalid rationale (a ``vertical 9:16'' aspect ratio misread as a $9.0$-versus-$16.0$-second duration).

\begin{figure}[tp]
  \centering
  \includegraphics[width=\linewidth]{fig_critic_cases.pdf}
  \caption{Six flaws identified through specialist measurement, one per row. Each row gives the
  specialist that answered, a five-frame filmstrip sampled uniformly across the clip, the human
  annotator's sentence verbatim, the specialist's own required-versus-measured pair with
  its evidence, and the resulting claim state. Displacements are normalized by the
  subject's radius, so no scene calibration is required. In the text case, the boxed
  banner region of the last frame is magnified in place, showing the render's own glyphs.
  Measurement text is copied from the run's records.}
  \label{fig:criticcases}
\end{figure}

\paragraph{What each measurement discriminates.} Each catches what an impression would not
(\cref{fig:criticcases}). The \emph{dolphin} leap is contradicted by a near-zero vertical excursion
a motion-only test would have passed; the \emph{bottle}'s surface and background move alike, so its
rolling is camera motion; the \emph{sun} is scored by peak departure (a rise then fall cancels) and
falls short amid comparable camera motion; \emph{counting} compares asserted against measured counts
exactly; and the \emph{text} head returns garbled glyphs that fail string equality, which a
generative model would silently normalize to the requested phrase.

\paragraph{Independent measurements and coverage.} A second dolphin sentence asserts the animal never
disappears, so two specialists evaluated the prompt's requirement that it disappear: the
vertical-travel check required downward motion but measured $2.07$ radii \emph{upward} ($-317$\,px
over $37$ frames, negative denoting up), and the persistence check required it to vanish but found it
on $121$ of $121$ frames with no gap of three or more; both independently contradict the
requirement. Each measurement also records its temporal coverage: in the banner case the recognition
head read $25$ of the clip's $241$ frames at stride ten, and a transcript stable across those $25$
evidences a persistent banner, whereas a briefly-appearing caption could fall between samples.

\paragraph{Specialist abstention.} In the run of \cref{tab:headline}, physical specialists were
called on $227$ checks, abstaining on $166$ and contradicting on $61$, each abstention carrying its
reason. The most common is no candidate grounding passing the semantic guard---shown one frame with
the candidate boxes and asked only whether they are the named subject: for seagulls scattering
behind fighter jets, the ``seagulls'' boxes held the jets and were discarded, yielding \unknownv{}
and separating a failure to find the subject from a violation by it.
\subsubsection{Agreement with human clip-level ratings}
\label{sec:vs2}
Since localized recall and clip-level scalars are not directly comparable, we compare \PhyReAct{} with VideoScore2~\cite{videoscore2_2025} on human clip-level semantic-adherence ratings, both on the $1$--$5$ scale and scored with VideoScore2's published script. The sample is a prespecified fixed-seed draw of $200$ VideoPhy-2~\cite{videophy2_2025} test clips ($55$ \emph{hard}-flagged) to which the scoring function was not fitted; agreement uses the four \cref{tab:vs2} columns, dropping per metric any unparsed prediction.

\begin{table}[t]
\centering\small
\renewcommand{\arraystretch}{1.2}
\begin{tabularx}{\textwidth}{@{}L c c c c c c@{}}
\toprule
 & \multicolumn{4}{c}{\textbf{All clips ($n=200$)}} & \multicolumn{2}{c}{\textbf{Hard subset ($n=55$)}} \\
\cmidrule(lr){2-5}\cmidrule(lr){6-7}
\textbf{Judge} & \textbf{Exact} & \textbf{Within 1} & \textbf{PLCC} & \textbf{SRCC}
 & \textbf{PLCC} & \textbf{SRCC} \\
\midrule
VideoScore2~\cite{videoscore2_2025} & \textbf{33.0\%} & \textbf{83.0\%} & 0.353 & 0.327 & 0.239 & 0.151 \\
\PhyReAct{} & 31.0\% & 76.0\% & \textbf{0.464} & \textbf{0.465} & \textbf{0.480} & \textbf{0.460} \\
\bottomrule
\end{tabularx}
\caption{Semantic-adherence agreement with human ratings on the $200$-clip VideoPhy-2
sample described in the text, scored with VideoScore2's evaluation script. The
physical-commonsense dimension is excluded because the critic does not emit a score on
that scale.}
\label{tab:vs2}
\end{table}

The critic trails on exact and within-one agreement but leads on PLCC and SRCC (\cref{tab:vs2}). A paired bootstrap gives SRCC gain $+0.138$ ($95\%$ CI $[+0.005,+0.277]$), PLCC gain $+0.110$ ($[-0.018,+0.242]$); on the $55$ hard clips, $+0.242$ ($[+0.018,+0.511]$) and $+0.309$ ($[+0.050,+0.583]$). This fits per-claim evidence degrading more gracefully than a single learned scalar on hard clips, though that subset is small and intervals wide. Caveats: the single $200$-clip draw describes that sample, not the benchmark; and the judges are not cost-matched (decomposition plus tool calls versus one forward pass), so this says nothing about efficiency.
\subsubsection{Held-out evaluation on Physion-Eval}\label{sec:physion}
Physion-Eval (\cref{sec:bench-vs}) holds out clips from this system's development: $9{,}569$
generated clips from five generators, each with a human report of physical glitches ($8{,}745$
with at least one). Two release limits bound what is measurable: the matched real-world reference
videos for the benchmark's $J$ statistic are withheld for copyright, and glitch-free clips are not
verified negatives. The release thus supports neither $J$ nor false-alarm rates for any
evaluator---only flag rates on release clips and matched-report rates on reported-glitch clips.

\PhyReAct{} runs over a fixed randomized stream; figures are interim, covering the first
$601$ clips ($543$ with a reported glitch, $58$ without). The in-context lesson state
(\cref{sec:icl-exp}) was active and updated as the stream progressed, so the run is prequential over
an evolving configuration, with no component seeing a label before scoring its clip.

Three measures separate prompt-grounded findings from physics-grounded reports (\cref{tab:physion}).
The \emph{matched-report clip rate} is the share of reported-glitch clips with a finding matching the
report under \cref{sec:matching}. The \emph{evidence-conditioned flag rate} is the share the served
model calls physically unrealistic from the check records alone---question, verdict, reason, and
numerical evidence---without frames or the reported defect. The \emph{anomaly-phrased finding rate}
is the share whose finding explicitly alleges a physical anomaly rather than an unmet prompt
requirement.

\Cref{fig:physionlb} places this work's evidence-conditioned flag rate against the
benchmark's published flag rate for every critic it evaluated, with the untrained-human
reference for scale; the full table, adding this work's other two rates and the
benchmark's $J$ statistics, is \cref{tab:physion} in \cref{sec:app-physion}. Because this
work's rows and the benchmark's rows use different samples and system configurations, their
juxtaposition is descriptive rather than a common-protocol comparison.

\begin{figure}[t]
  \centering
  \includegraphics[width=\linewidth]{fig_physion_leaderboard.pdf}
  \caption{Physion-Eval physical-implausibility flag rate for every evaluator of
  \cref{tab:physion}, sorted high to low, with the untrained-human reference (light) and
  this work's evidence-conditioned rate (red). This work's rate is measured on the first
  $601$ release clips; the benchmark's critic rates are transcribed from its Figure~4 over
  its full evaluation set, so the comparison is descriptive rather than common-protocol.}
  \label{fig:physionlb}
\end{figure}

The evidence-conditioned flag rate is nearly twice the anomaly-phrased finding rate,
indicating that the accumulated evidence supports physical-implausibility
judgements beyond what individual findings' wording states.
\subsection{Context-training self-evolution of the critic}
\label{sec:icl-exp}
The inward channel distills lessons from misses into the critic's readable state, improving without touching a weight; the formalism---distillation map \eqref{eq:distill}, acceptance gate \eqref{eq:elitism}, edited state $S=(\mathcal{O},M,K)$---is in \cref{sec:framework}. We measure flaw recall.

On a plan's miss of a human-annotated flaw, $\operatorname{Distill}$ \eqref{eq:distill} reads the trace and miss and writes a lesson naming a claim property and the action it licenses (e.g.\ turning a singular-entity or exact-quantity requirement into a \emph{separate cardinality check verified across beginning, middle and end}); curation groups candidates into $33$ \emph{families} of equivalent advice, one lesson each, appended to the planner's rubric. Only $M$ changes ($S' = (\mathcal{O}, M', K)$); a family's name labels the defect its lesson concerns.

Implementation departs from the formalism twice. First, all $33$ lessons go on every clip: the selective form (one relevance call per lesson--claim pair) costs $132$ model calls for a four-claim prompt and $495$ for a fifteen-claim one, whereas all-at-once costs no calls and adds $9$\,kB to a planner prompt already $17.8$\,kB long. Second, the acceptance gate \eqref{eq:elitism} is dropped: $U_{\mathrm{val}}$ scores only rules already seen in a run, not one from a fresh miss---diagnostically, on $80$ earlier candidates over a ten-clip set it admitted none ($39$ failed its minimum-evidence criterion, $41$ never triggered). The evaluated intervention is the curated $33$ lessons, without the gate.

\paragraph{Protocol and result.} Lessons distilled from experience clips are tested on held-out ones: excluding every clip sharing an identifier or restating an experience-clip prompt leaves $306$ clips with $502$ flaws, so no lesson is evaluated on the prompt it was written from. Every condition scores those flaws under \cref{sec:matching}'s protocol (a flaw is \emph{found} when the critic covers it wholly or in part), untaught loading an empty lesson file through the same path; five nested conditions---no lessons, then $8$, $13$, $19$ and $33$ lessons from $25$, $50$, $75$ and $117$ clips, each a prefix of the next---form a dose curve. With all $33$ the planner finds $\mathbf{375}$ of $502$ flaws ($\mathbf{74.7\%}$) against $340$ untaught ($67.7\%$): $62$ switch missed-to-found and $27$ back, net $+35$ at $p = 3\!\times\!10^{-4}$ (paired sign test). More experience does not extend it (\cref{fig:iclaggregate}): the four taught conditions find $375$, $377$, $374$ and $375$, each $+34$ to $+37$ over untaught at $p < 10^{-3}$ but no two differing by more than $3$ ($p \geq 0.8$ pairwise), and the unfiltered set ($363$ clips, $626$ flaws) matches. Later lessons reshuffle which flaws are found, not how many---$36$ move each way between the $8$- and $33$-lesson conditions---so the gain is present already with lessons from $25$ clips.

\begin{figure}[t]
  \centering
  \begin{minipage}[c]{0.47\textwidth}
    \centering
    \includegraphics[width=\linewidth]{fig_icl_aggregate.pdf}
  \end{minipage}\hfill
  \begin{minipage}[c]{0.49\textwidth}
    \caption{Flaws found against the number of experience clips the lessons were distilled from,
    over the same $502$ held-out flaws in every condition. The dotted line marks the untaught
    condition, $340$ of $502$ ($67.7\%$).
    Whiskers are $95\%$ intervals on the change from untaught, in flaws, computed from the flaws that
    switch status, since the conditions score the same flaws. Every taught condition exceeds the
    untaught one, while pairwise comparisons among the taught conditions return $p \geq 0.8$.}
    \label{fig:iclaggregate}
  \end{minipage}
\end{figure}

\paragraph{Effect by flaw category.} Both conditions score the same flaws, so we report the paired change and its interval from flaws that switched status (\cref{tab:iclcategory}). Between untaught and $33$-lesson conditions, two categories' $95\%$ interval excludes zero: \emph{action or event occurs}, largest at $211$ flaws, gains $17$ ($8.1$ percentage points, $[+2, +12]$); \emph{on-screen text} gains $8$ of $73$ ($11.0$ points, $[+2, +15]$). All other intervals contain zero, and reading smaller movements as rates overstates them: $+7$ points for camera and style is one flaw of $14$, $+16$ for counting is three of $19$. Flaws rise in both severity grades holding the bulk of the set---$19$ of $288$ at severity~$3$ ($p = 0.01$) and $13$ of $163$ at severity~$4$ ($p = 0.03$)---the remaining $51$ at grades $1$, $2$ and $5$.

% GENERATED by techreport/figs/emit_icl_category_tex.py -- do not edit by hand.
% Paired comparison: the three conditions score the same flaws.

\begin{table}[t]
\centering\small
\setlength{\tabcolsep}{5pt}\renewcommand{\arraystretch}{1.15}
\begin{tabular}{@{}p{3.9cm} c c c c c c c@{}}
\toprule
 & & \multicolumn{3}{c}{\textbf{Flaws found, with}} & & & \\
\cmidrule(lr){3-5}
\textbf{Flaw category} & \textbf{Flaws} & \textbf{no lessons} & \textbf{8} & \textbf{33} & \textbf{Change} & \textbf{95\% interval} & \textbf{Lessons} \\
\midrule
\textbf{order and timing} & 15 & 6 & 7 & 10 & +4 (+27) & $[+1, +27]$ & 3 \\
how many & 15 & 8 & 9 & 10 & +2 (+13) & $[-11, +24]$ & 2 \\
on-screen text & 55 & 36 & 46 & 43 & +7 (+13) & $[-0, +20]$ & 1 \\
\textbf{action or event occurs} & 177 & 124 & 139 & 143 & +19 (+11) & $[+5, +14]$ & 13 \\
camera and style & 13 & 11 & 12 & 12 & +1 (+8) & $[-5, +8]$ & 0 \\
attribute & 38 & 24 & 24 & 26 & +2 (+5) & $[-4, +10]$ & 2 \\
object identity & 42 & 33 & 36 & 34 & +1 (+2) & $[-1, +2]$ & 3 \\
physical motion & 12 & 8 & 7 & 8 & +0 (+0) & $[-14, +14]$ & 0 \\
spatial relation & 40 & 25 & 24 & 23 & -2 (-5) & $[-12, +6]$ & 1 \\
\bottomrule
\end{tabular}
\caption{Effect of the distilled lessons on each flaw category. A flaw is \emph{found} when the critic's findings cover it wholly or in part, as in \cref{tab:headline}. The three middle columns count that category's flaws found under no lessons, under $8$ lessons distilled from $25$ experience clips, and under $33$ from $117$; the change is given in flaws with percentage points in brackets, since several categories carry few enough flaws that a rate alone would overstate the movement. Because the three conditions score the same flaws the comparison is paired, and the interval on the change is computed from the flaws that switched status. \textbf{Bold} marks the two categories whose interval excludes zero. The last column counts the lesson families added between $25$ and $117$ experience clips whose name refers to that category.}
\label{tab:iclcategory}
\end{table}


\paragraph{Where the lessons point.} Gains are not aligned with lesson names. All eight families from the first $25$ clips concern actions, none mentioning text, counting or spatial relations, yet at $8$ lessons \emph{on-screen text} posts the largest percentage-point gain of any category ($10$ of $73$) before any lesson mentions text, then loses two after one does; the later $25$ families do not move their nominal targets. This fits the first lessons supplying a general instruction---separate a compound requirement into checkable parts---applying well beyond the categories they were written from, which the plan measurements below assess.

\paragraph{Effect of the lessons on the plan.} Lessons compile the prompt into more claims, not more checks per claim: from untaught to $33$ lessons, checks per claim rise only $1.13\!\to\!1.18$ while typed physical claims rise $2{,}905\!\to\!5{,}106$ over $306$ clips ($9.5\!\to\!16.7$ per clip)---as the decomposition reading predicts. This locates the flattening: claims per clip keep rising across all five conditions ($9.5$, $15.1$, $16.0$, $16.1$, $16.7$), so later lessons keep changing the plan while flaws found stall. On-screen text shows this at clip level: text \emph{checks} rise $81\!\to\!136$ (already $114$ at $8$ lessons), yet of $56$ text-flaw clips the number getting any text check does not ($41$ untaught vs.\ $40$ taught)---lessons add checks without widening the set of clips checked.
\subsection{Writing a verdict back into generation}\label{sec:refine}
The critic's outward channel---a verdict with its evidence, fed back as a prompt rewrite that
regenerates the clip---is claimed but not yet measured. We close that loop on a held-out generator
and ask a deliberately narrow question: does it help on the metric that triggered it, and does an
\emph{independent} rater register the same change?

The loop composes the two maps of \cref{eq:gen-crit-maps} through a rewrite operator: from a base
clip $\hat{\mathcal{V}}\sim\operatorname{Gen}(p,C)$ the critic returns
$\operatorname{Crit}(\hat{\mathcal{V}},p)=(v,\pi,\mathcal{F})$---a verdict
$v\in\{\textsf{plausible},\textsf{implausible},\textsf{abstain}\}$, a plausibility score
$\pi=p_{\text{plausible}}$, and the violated obligations $\mathcal{F}=\{(e_j,a_j)\}_j$ with their
evidence. A text-only rewrite $\rho$ folds those failures into the prompt and
the \emph{same} generator resamples under the same control and seed,
\begin{equation}
  p' = \rho\bigl(p,\mathcal{F}\bigr)\ \ \text{when}\ \ v=\textsf{implausible}\ \ (\text{else}\ p'=p),
  \qquad
  \hat{\mathcal{V}}' \sim \operatorname{Gen}\bigl(p',C\bigr),
  \label{eq:refine-loop}
\end{equation}
so the channel acts only through the prompt---diffusion weights, control $C$, mask, and seed
untouched---changing appearance while fixing the motion path. We run one round, reading the change
in $\pi$ and $v$ against an independent rater applied to both.

\paragraph{Protocol.} We run the full \emph{evaluate\,$\to$\,rewrite\,$\to$\,regenerate} loop over
Cosmos3-Nano, a generator unused anywhere in this system's development, on its $600$ VideoPhy-2 test
clips ($589$ complete the loop with a paired score, so the before/after columns rest on $600$ and
$589$). Each baseline and its regeneration are scored two ways: by \PhyReAct{}'s own metric
($p_{\text{plausible}}$ and the verdict of \cref{sec:eval}), the quantity the rewrite optimizes; and
by VideoPhy-2's official AutoRater~\cite{videophy2_2025}, an \emph{independent} rater returning
semantic adherence (SA), physical commonsense (PC), and their joint agreement (both $\ge 4$). The
AutoRater is the control---not the loop's objective and sharing no component with the critic.

\paragraph{On its own metric the loop is effective.} By \PhyReAct{}'s own measurement
(\cref{tab:refine-selfmetric}) the loop is targeted with low collateral: it acts on the $14.8\%$ of
clips it judges flawed, raises $p_{\text{plausible}}$ on $72.4\%$ of those and flips $49$ of $65$
implausible verdicts to plausible, regressing only $1.6\%$ of unflagged clips.

% PhyReAct's own physical metric (p_plausible / verdict) over the 589-clip loop.
% These are the critic's SELF-assessed quantities: the same critic both triggers a
% rewrite and scores the result, so this table is not an independent measurement.
% Numbers supplied from critic_full.jsonl; not recomputed here.

\begin{table}[t]
\centering\small
\renewcommand{\arraystretch}{1.2}
\begin{tabular}{@{}p{9.6cm} r@{}}
\toprule
\textbf{Quantity (PhyReAct self-assessed)} & \textbf{Value} \\
\midrule
Clips PhyReAct judged flawed and rewrote (\emph{acted})   & $87/589$ \;($14.8\%$) \\
\quad of these, $p_{\text{plausible}}$ increased          & $63/87$ \;($72.4\%$) \\
\quad verdict flipped \emph{implausible}$\to$\emph{plausible} & $49/65$ \;($75\%$) \\
\quad mean $\Delta p_{\text{plausible}}$ (acted)          & $+0.524$ \\
Collateral damage: non-flagged clips scored worse         & $8/502$ \;($1.6\%$) \\
Overall $p_{\text{plausible}}$, baseline $\to$ rewrite     & $0.797 \to 0.862$ \\
\bottomrule
\end{tabular}
\caption{The same loop scored by \PhyReAct{}'s own physical metric. Every quantity
here is self-assessed: the critic that triggers a rewrite is also the judge of its
result, so these figures show the loop is internally consistent---targeted (it acts on
$14.8\%$ of clips), effective on its targets ($72.4\%$ improve, $49$ flip from
implausible to plausible), and low-collateral ($1.6\%$)---but they are not an
independent confirmation. They should be read against \cref{tab:refine-autorater},
where an independent rater registers no corresponding change, including on the
physical dimension. The $65$ denominator is the acted clips PhyReAct had first judged
implausible.}
\label{tab:refine-selfmetric}
\end{table}

% VideoPhy-2 official AutoRater (mplug-based), baseline generation vs the
% VeriPhy refinement loop, over all 589 rewritten clips. Numbers supplied from
% the run's joint/scores CSVs; not recomputed here.

\begin{table}[t]
\centering\small
\renewcommand{\arraystretch}{1.2}
\begin{tabularx}{\textwidth}{@{}L cc cc@{}}
\toprule
 & \multicolumn{2}{c}{\textbf{All ($n=589$)}} & \multicolumn{2}{c}{\textbf{Hard subset}} \\
\cmidrule(lr){2-3}\cmidrule(lr){4-5}
\textbf{VideoPhy-2 AutoRater metric} & \textbf{Baseline} & \textbf{Rewrite}
 & \textbf{Baseline} & \textbf{Rewrite} \\
\midrule
Mean semantic adherence (SA)      & 3.12   & 3.13   & 2.75  & 2.76  \\
Mean physical commonsense (PC)    & 3.59   & 3.58   & 3.36  & 3.32  \\
SA $\ge 4$                        & 27.8\% & 28.0\% & 8.9\% & 7.9\% \\
PC $\ge 4$                        & 55.0\% & 54.5\% & 37.2\% & 36.2\% \\
Joint (SA $\ge 4$ \emph{and} PC $\ge 4$) & 22.8\% & 23.1\% & 3.9\% & 3.4\% \\
\bottomrule
\end{tabularx}
\caption{The refinement loop scored by an \emph{independent} rater---VideoPhy-2's
own AutoRater~\cite{videophy2_2025}---on the same $589$ Cosmos3-Nano clips before and
after the loop. No cell moves beyond the rater's noise: joint agreement shifts by
$+0.3$ points over all clips and $-0.5$ on the hard subset, and both underlying
dimensions are flat to slightly lower, including the physical dimension the loop is
meant to improve (mean PC $3.59\!\to\!3.58$ all, $3.36\!\to\!3.32$ hard). On the hard
subset every entry is unchanged or slightly lower. The rewrite therefore produces no
change this independent rater can register, in either the semantic or the physical
dimension. Hard-flagged clips follow VideoPhy-2's designation.}
\label{tab:refine-autorater}
\end{table}


\begin{figure}[t]
  \centering
  \includegraphics[width=\linewidth]{fig_refine_data.pdf}
  \caption{The refinement loop read two ways, on the same $589$ Cosmos3-Nano clips.
  \textbf{Left:} by \PhyReAct{}'s own physical metric---the quantity the rewrite
  optimizes---the loop is targeted and low-collateral: of the clips it flagged, $63/87$
  improve and $49/65$ flip from implausible to plausible, against $8/502$ unflagged
  clips that regress. \textbf{Right:} the same clips scored by VideoPhy-2's independent
  AutoRater show no movement---each dumbbell from baseline (navy) to rewrite (red) collapses
  to a point at every threshold, on both the full set and the hard subset, in the semantic
  \emph{and} the physical dimension. Every number is transcribed from
  \cref{tab:refine-selfmetric,tab:refine-autorater}. The improvement is thus
  visible on the critic's own metric and not on the one rater independent of it. Clips in
  \cref{fig:refinefilm}.}
  \label{fig:refinedata}
\end{figure}

\begin{figure*}[tbp]
  \centering
  \includegraphics[width=\linewidth]{fig_refine_film.pdf}
  \caption{What the critic flagged, and the clip the rewrite regenerated---clips
  \PhyReAct{} judged implausible and resampled through \cref{eq:refine-loop}. Each case
  is two real filmstrips: the \emph{baseline} clip (grey) and the clip regenerated from
  the rewritten prompt (\emph{rewrite}, red), six frames early to late; red boxes mark
  the cited physical violation and its repair, and the full regenerated prompt is printed
  verbatim beneath. Cases were chosen by an independent frame-by-frame review under an
  honesty constraint---the before/after must be the \emph{same} scene, not a rewrite that
  swaps the scene or introduces a new artifact---which most flagged clips fail; these are
  cases where the fix is visible and the comparison fair. Three further cases are in
  \cref{fig:refinefilm2,fig:refinefilm3}; the quantitative reading is \cref{fig:refinedata}.}
  \label{fig:refinefilm}
\end{figure*}

\begin{figure*}[tbp]
  \centering
  \includegraphics[width=\linewidth]{fig_refine_film2.pdf}
  \caption{Critic-guided rewrite, continued, in the layout of \cref{fig:refinefilm}:
  coins that fuse into a shapeless blob and a frisbee that warps into a red ball
  mid-flight, each held to a single rigid object by the rewrite. Red boxes mark the
  flagged defect and its repair; the full rewritten prompt is printed verbatim.}
  \label{fig:refinefilm2}
\end{figure*}

\begin{figure*}[tbp]
  \centering
  \includegraphics[width=\linewidth]{fig_refine_film3.pdf}
  \caption{Critic-guided rewrite, continued, in the layout of \cref{fig:refinefilm}: a
  soccer ball that duplicates during the kick---two balls in the baseline (boxed), one
  in the rewrite. The full rewritten prompt is printed verbatim.}
  \label{fig:refinefilm3}
\end{figure*}

\paragraph{An independent rater registers no change.} The same $589$ clips scored by VideoPhy-2's
AutoRater show no movement in either dimension (\cref{tab:refine-autorater}; \cref{fig:refinedata}
places the two readings side by side, clips in \cref{fig:refinefilm}): joint agreement, SA, and---%
critically---the \emph{physical} PC dimension all stay flat to lower, within rater noise. Stated
together, not collapsed: the gain appears only on the metric the critic itself assigns, while the
one rater independent of the loop registers no change on either axis. The design cannot separate two
readings---that rewrites improve a physical quality the critic can measure but the AutoRater cannot,
or that they raise the critic's own score without improving the clip---because the critic both
triggers and judges the effect it reports. (External position only: Cosmos3-Nano's AutoRater joint
score is unchanged by the rewrite, $22.8\!\to\!23.1\%$ (all) and $3.9\!\to\!3.4\%$ (hard), near
published human-rated leaderboard entries; but human and AutoRater joint scores share no common
protocol, so this is descriptive only.) Every reported figure was recomputed independently from the
run's per-clip files and reproduces exactly (\cref{tab:refine-verify} in \cref{sec:app-refine-verify}).

\paragraph{Scope.} Confirming the rewrites improve physical fidelity as an independent judge would
score it needs a rater both independent of the critic and sensitive to the targeted physical
property---which the AutoRater's flat PC suggests it is not on these clips. All $589$ per-clip
comparisons and both scoring streams ship with the run.
\subsection{Scope and interpretation}\label{sec:eval-scope}
Each setting must be read on its own terms. The generator analyses (\cref{sec:gen-experiments}) are
pre-registered probes on small closed-form scenes, characterizing the deployed pipeline rather than
its generality. Core recall uses recorded annotations on development data, where an unrecorded defect
cannot be scored, and its question-decomposition baseline was specified after inspecting critic
outputs. The context-training result (\cref{sec:icl-exp}) is held out from the distillation pool but
still drawn from the development corpus. Only the interim, prequential Physion-Eval run
(\cref{sec:physion}) is held out from development entirely---with an evolving lesson state, a
different annotation vocabulary, and no frozen-core protocol. The refinement loop
(\cref{sec:refine}) is closed on a held-out generator but confirmed only on the critic's own metric;
the one independent rater shows no change in either dimension.
